\documentclass[11pt,a4paper]{article}

% ---------------------------------------------------------------------------
%  Evra_cost_benefit.tex
%  Working draft: cost-benefit analysis of Iceland adopting the euro (ESB/EU).
%  This is a source/draft document. It does NOT publish on kannski-esb.is,
%  which is an HTML-only static site. Content can be pasted/edited below.
% ---------------------------------------------------------------------------

\usepackage[utf8]{inputenc}
\usepackage[T1]{fontenc}
\usepackage{lmodern}
\usepackage[icelandic,english]{babel}
\usepackage{geometry}
\geometry{margin=2.5cm}
\usepackage{booktabs}
\usepackage[authoryear,round]{natbib}
\usepackage{hyperref}
\hypersetup{colorlinks=true, linkcolor=blue, urlcolor=blue, citecolor=blue}

\title{Evra: Cost--Benefit Analysis \\ \large Kostnaður og ábati við upptöku evru}
\author{Kannski --- kannski-esb.is}
\date{\today}

\begin{document}
\maketitle

\begin{abstract}
This note weighs the main costs and benefits for Iceland of adopting the euro.
The central benefit is a durable fall in real interest rates --- on the order
of 3--4 percentage points for household borrowing, the range reported in
studies of earlier adopters, with Spain, Portugal, Italy and the Baltic states
the leading examples --- which would act as a large transfer of income from
banks to households, and as cheaper credit for firms. The central cost is the
loss of an independent monetary policy. Two developments have shrunk that
cost: labour mobility in and out of Iceland is far higher than in the 1990s,
and, as the 2026 expert report for the Ministry of Finance led by Catherine
L.~Mann documents, Iceland's business cycle is now moderately synchronised
with the euro area --- output correlations around 0.7 since 2009 --- while the
kr\'ona has tended to amplify rather than absorb shocks. The report concludes
that the cost of the kr\'ona is substantial and likely exceeds, on average,
the benefit of an independent currency.
\end{abstract}

\section{Introduction}
The question is not whether the euro has costs --- it plainly does --- but
whether, for Iceland today, the benefits outweigh them. The honest case for
adoption rests on one large, measurable gain and concedes one real, structural
loss. The gain is lower real interest rates; the loss is the surrender of a
sovereign monetary policy. Opponents of EU membership keep some credibility
only by conceding the first and pressing hard on the second. This note sets
out both sides and explains why the author now gives the main counter-argument
less weight than he did in a 1997 B.A.\ thesis on the subject. It also draws
on the expert report on Iceland's currency options commissioned by the
Ministry of Finance and Economic Affairs and delivered in May 2026 by a group
of four international economists chaired by Catherine L.~Mann, former chief
economist of the OECD and member of the Bank of England's Monetary Policy
Committee \citep[see][]{mann2026}.

\section{Why Icelandic interest rates are so high}
Any cost--benefit account has to start here, because the size of the benefit is
just the size of the gap that adoption would close. Real mortgage rates in
Iceland sit at roughly 4.5--5.5\,\% depending on loan form, far above the euro
area. Four forces keep them there.

\begin{enumerate}
  \item \textbf{The ``Iceland premium'' (the peso problem).} Rates are
  generally higher in countries with an unstable currency, to compensate
  holders for the risk that the currency falls. If the exchange rate has in
  fact been stable and nominal rates are \emph{still} much higher than in the
  euro area, that gap is precisely the ``peso problem'', named for Mexico's
  experience: the market is pricing a devaluation that has not yet happened
  but is believed possible.
  \item \textbf{Monetary-policy risk premium.} Chronic instability and a
  low-credibility monetary regime leave the future path of inflation uncertain,
  and that uncertainty is itself priced as a risk premium.
  \item \textbf{Small currency, small market $\rightarrow$ oligopoly.} The
  króna is a tiny currency and the market is small. It rarely pays a foreign
  financial firm to keep specialists forecasting Icelandic exchange rates,
  inflation and growth, plus specialists in Icelandic law, merely to compete for
  Icelandic mortgages. The result is a domestic banking oligopoly (fákeppni)
  under little pressure to offer internationally competitive rates.
  \item \textbf{Post-2008 regulation.} Regulation of Icelandic banks has
  plausibly become quite stiff since the 2008 crash, making them less efficient
  than comparable banks abroad --- and with no foreign entrants on the market,
  there is little pressure to relax it.
\end{enumerate}

The Mann report corroborates the risk-premium channels: high-frequency
exchange-rate volatility, punctuated by episodic large moves, feeds risk
premia on domestic interest rates; recurrent depreciations sustain them; and
the resulting high domestic rates hamper growth in interest-sensitive sectors
and hurt indebted households and firms. It also notes that Iceland's gross and
net capital flows have been three to four times more volatile, as a share of
GDP, than those of comparable countries over the last 25 years
\citep{mann2026}.

\section{Benefits}

\subsection{Lower real interest rates}
The experience of countries that adopted the euro --- Spain, Portugal, Italy and
the Baltic states --- is that banks in neighbouring countries began competing
for ordinary lending against domestic banks, and real rates fell; in the
Baltic case this began already in the run-up to adoption. The general
conclusion of the studies the author has seen puts the reduction, measured in
real rates, at 3--4 percentage points, with even larger falls in some
countries, Portugal for example. He sees no reason to expect a different
outcome in Iceland. Crucially, this development came in direct connection with
euro adoption; it is hard to see how the same competitive entry could be
achieved without such a step.

Two figures should be kept apart here. The Mann report's cross-country
estimates concern the economy-wide cost of capital, which fell by 25--100
basis points in earlier adopters, with sovereign rates converging to euro-area
levels over time \citep{mann2026}. The 3--4 point figure concerns
\emph{household} borrowing in the 1990s convergence episodes, where starting
spreads were far wider --- as Iceland's are today. The two are consistent:
the bigger the initial wedge, and the weaker the initial banking competition,
the bigger the fall.

\subsection{Transaction costs and access to capital}
Beyond interest rates, the report puts numbers on the mechanical gains. For
Iceland, where the bid-ask spread on the kr\'ona is very high and formal
foreign-exchange trading is thin, eliminating currency-conversion costs is
worth a recurring saving of roughly 0.25\,\% of GDP (Croatia saved an
estimated 0.19\,\%, Bulgaria and Lithuania around 0.6\,\%). External
financing, public and private, would be cheaper, and the state would gain a
stronger presence on international bond markets \citep{mann2026}.

\subsection{Household savings and equity build-up}
The saving to households is the reduction in the real rate they pay the banks,
multiplied by the size of their loan portfolio. Take the calculator's baseline:
a new mortgage of about 60\,m.kr.\ (85\,\% of a 70.6\,m.kr.\ first-buyer home,
HMS 2025) over 25 years. Lowering the real rate by several points transforms the
composition of every payment.

It is true, as critics note, that house prices rise when rates fall --- and one
should expect that. But two things follow. First, cheaper credit creates an
incentive to build more housing: investment in new housing in Spain roughly
doubled at euro adoption. Second, even if the monthly payment stayed similar
but at a lower real rate (a premise the author considers out of line with the
experience of most countries), far more of it would be genuine \emph{saving}
in the household's own equity --- the house, now worth more --- rather than
pure interest cost that embeds no saving at all. Later in life, when people
downsize, they would hold a far more valuable asset to sell.

\subsection{Investment and growth}
Over the longer run, lower real rates amount to a direct transfer of income from
banks --- which charge higher real rates than competitors abroad --- to
Icelandic households. That transfer would be an enormous windfall. Lower rates
would also raise the incentive to invest, the foundation of growth, not least
for the small and medium-sized firms that today cannot borrow in euros at all.

\section{The Mann report: Iceland's cycle and the European cycle}

\subsection{What was asked and what was answered}
In 2026 the Ministry of Finance and Economic Affairs received the report it
had commissioned on Iceland's currency options, written by Catherine L.~Mann
(chair), Hilde C.~Bj{\o}rnland (BI Norwegian Business School), Steinar Holden
(University of Oslo) and Stijn Claessens (formerly of the Bank for
International Settlements, now at Yale). Its mandate was to weigh an
independent currency against membership of a larger currency area such as the
euro zone. Its headline finding: the cost of the kr\'ona is substantial and
likely exceeds, on average, the benefit of an independent currency; the
exchange rate has not fostered stability except in major crises, and its
swings have often amplified rather than dampened economic shocks
\citep{mann2026}.\footnote{The report frames itself as groundwork for informed
debate rather than a recommendation to adopt, and the government's summary of
it has been politically contested in Iceland.}

\subsection{Synchronisation with the European cycle}
The classic objection to euro membership --- the one this author leaned on in
1997 --- assumes that Iceland's cycle diverges from Europe's, so that ECB
policy would systematically misfit Icelandic conditions. The report's evidence
undercuts the premise. Iceland's output-growth correlation with the euro area
is about 0.65 over the full sample and 0.71 after 2009 (output-gap
correlation about 0.70); correlations with the individual Nordic countries
run at 0.6--0.65 post-2009. The major turning points --- the 2008--09 crisis,
the pandemic, the recovery --- were plainly shared. The report's own words:
business-cycle synchronisation ``is not a binding obstacle to euro adoption,
and may potentially rise further over time through the deeper trade and
financial integration that euro membership itself encourages'', citing the
endogenous currency-area result of \citet{frankelrose1998}. In other words,
synchronisation is not only already moderate-to-high; adopting the euro would
raise it further.

\subsection{The kr\'ona: shock absorber or shock source?}
The textbook defence of a float is that the exchange rate absorbs shocks. The
report finds Iceland's experience does not cleanly fit: output and inflation
volatility remain the highest in the Nordic group even after 2009, and the
kr\'ona has moved largely independently of the other Nordic floating
currencies even though the underlying real economies co-move --- which is what
one would expect if the currency were driven by risk-premium and financial
shocks rather than by fundamentals. Depreciation episodes are associated with
higher consumer prices and a persistent compression of real wages, with no
evidence of sustained gains in export volumes. Meanwhile Iceland's policy rate
already moves closely with euro-area rates, so the practical scope of the
``independent'' policy being defended is narrower than the textbook suggests.
Fixing to the euro would itself remove a source of macroeconomic noise that
has historically passed straight through to inflation and output
\citep{mann2026}.

\section{Costs}

\subsection{Loss of independent monetary policy}
The real cost is that euro adoption removes Iceland's capacity to run a
flexible monetary policy suited to domestic conditions --- above all, to let
the currency fall in a downturn, cutting real wages and thereby unemployment.
In Mundell's theory of optimum currency areas --- the work for which he later
won the Nobel Prize \citep{mundell1961} --- the cost of a shared currency is
greatest when labour does not move easily between regions hit by different
shocks and when nominal wages are sticky. Icelandic nominal wages probably
would not fall much in a downturn, and little may have changed there since
1997. The Mann report adds a genuine caveat on the other side of its
synchronisation finding: the \emph{amplitude} of Iceland's fluctuations is
roughly twice that of Denmark or Norway, and Iceland has had idiosyncratic
episodes --- the 2004--08 boom-bust, the 2015--18 tourism expansion --- with
no close Nordic counterpart. Giving up the domestic instrument would
``not be costless'' \citep{mann2026}.

\subsection{Why the cost has shrunk: labour mobility}
Against this, labour mobility into and out of Iceland is now far greater than
it was in, say, 1997. Iceland today imports a great deal of labour from the
Schengen area, and the European internal market has come to resemble the
American one, where people move between regions in pursuit of work. A large
share of those working in Iceland do so temporarily --- many hard-working
Poles, for instance, come to earn income for their families and then return
home. In a downturn, arrivals from the European Economic Area would fall
sharply and the flow would reverse into outflow. That outflow cushions
unemployment and so \emph{reduces} the cost of giving up an independent
monetary policy. This is the main reason the author now weighs the Mundell
objection less heavily than he did writing his first essay on the topic at his
1997 B.A.\ graduation from the University of Iceland --- and the Mann report's
synchronisation evidence pushes the same way: the more Iceland's cycle moves
with Europe's, the less often a separate policy would even be wanted.

\subsection{Preconditions and transition risks}
The report is explicit that adoption shifts the burden of stabilisation onto
fiscal policy and the labour market, and it judges labour-market reform
necessary under \emph{any} currency arrangement --- but a precondition for
successful membership of a monetary union. It flags two further points
deserving care: Iceland would be the first country to join with so high a GDP
per capita, which makes the entry exchange rate a key variable, and Iceland's
demographics imply a long-run neutral real rate that may differ from the
euro-area's \citep{mann2026}. Rising house prices, noted above, are a genuine
distributional effect, and the transition itself --- including possible
speculative pressure --- would need careful management. These belong on the
ledger; on the argument here they do not outweigh the interest-rate gain.

\section{Conclusion}
The euro's benefit for Iceland is concrete and quantifiable --- a lasting fall
in real interest rates that transfers income from banks to households and
lowers the cost of credit for firms, lifting investment, plus a recurring
transaction-cost saving of around a quarter of a percent of GDP. Its cost is
the loss of an independent monetary policy. That cost weighed heavily in the
1990s. It weighs less now, for two reasons that reinforce each other: labour
now moves in and out of Iceland at a scale that cushions downturns, and
Iceland's business cycle has become moderately synchronised with the euro
area's, while the kr\'ona itself has behaved more like a source of shocks than
a buffer against them. When the government's own commissioned experts conclude
that the kr\'ona's costs likely exceed its benefits, the honest answer has
moved, from the scepticism of 1997 towards, at the very least, \emph{kannski}.

\begin{thebibliography}{9}

\bibitem[Frankel and Rose(1998)]{frankelrose1998}
Frankel, J.~A. and A.~K. Rose (1998). ``The Endogeneity of the Optimum
Currency Area Criteria.'' \emph{The Economic Journal} 108(449), 1009--1025.

\bibitem[Mann et al.(2026)]{mann2026}
Mann, C.~L., H.~C. Bj{\o}rnland, S.~Holden and S.~Claessens (2026).
\emph{Currency Matters: A Report for the Iceland Ministry of Finance and
Economic Affairs} (Sk\'yrsla um valkosti \'Islands \'i gjaldmi\dh lam\'alum).
Reykjav\'ik, May 2026.
\url{https://www.stjornarradid.is/library/02-Rit--skyrslur-og-skrar/Sk%C3%BDrsla%20um%20valkosti%20%C3%8Dslands%20%C3%AD%20gjaldmi%C3%B0lam%C3%A1lum%20-%20001.pdf}

\bibitem[Mundell(1961)]{mundell1961}
Mundell, R.~A. (1961). ``A Theory of Optimum Currency Areas.''
\emph{American Economic Review} 51(4), 657--665.

\end{thebibliography}

\end{document}
